\chapter{Connection between Minkowski and Euclidean Field Theory and with Statistical Mechanics}

This section is a brief recapitulation of technical issues regarding the Euclidean path integral and extracting physics from Euclidean correlation functions. For details see [3,4] or the texts on statistical field theory, for example [22,23].

A $D$-dimensional Minkowski field theory ($D-1$ spatial dimensions and one time dimension) is connected to a $D$-dimensional Euclidean field theory through analytical continuation. Under Wick rotation
\begin{equation}
\begin{aligned}
x_0 &\equiv t \to -i x_4 \equiv -i\tau , \\
p_0 &\equiv E \to i p_4 .
\end{aligned}
\label{eq:5.1}
\end{equation}

The Euclidean convention is
\begin{equation}
\begin{aligned}
x_E^2 &= \sum_{i=1}^{4} x_i^2 = \mathbf{x}^2 - t^2 = -x_M^2 , \\
p_E^2 &= \sum_{i=1}^{4} p_i^2 = \mathbf{p}^2 - E^2 = -p_M^2 .
\end{aligned}
\label{eq:5.2}
\end{equation}

The connection between statistical mechanics and a Euclidean field theory in $D$-dimensions is most transparent when the field theory is quantized using the Feynman path integral approach. This connection forms the basis of numerical simulations using methods common to those used to study statistical mechanics systems. A synopsis of the equivalences between the two is given in Table~\ref{tab:1}.

\begin{table}[htbp]
\centering
\caption{The equivalences between a Euclidean field theory and Classical Statistical Mechanics.}
\label{tab:1}
\begin{tabular}{ll}
\toprule
\textbf{Euclidean Field Theory} & \textbf{Classical Statistical Mechanics} \\
\midrule
Action & Hamiltonian \\
unit of action $\hbar$ & units of energy $\beta = 1/kT$ \\
\shortstack{Feynman weight for amplitudes \\ $e^{-S/\hbar} = e^{-\int L\,dt/\hbar}$} & Boltzmann factor $e^{-\beta H}$ \\
\shortstack{Vacuum to vacuum amplitude \\ $\int D\phi\, e^{-S/\hbar}$} & \shortstack{Partition function \\ $\sum_{\mathrm{conf.}} e^{-\beta H}$} \\
Vacuum energy & Free Energy \\
Vacuum expectation value $\langle 0|O|0\rangle$ & Canonical ensemble average $\langle O\rangle$ \\
Time ordered products & Ordinary products \\
Green's functions $\langle 0|T[O_1 \ldots O_n]|0\rangle$ & Correlation functions $\langle O_1 \ldots O_n\rangle$ \\
Mass $M$ & correlation length $\xi = 1/M$ \\
Mass-gap & exponential decrease of correlation functions \\
Mass-less excitations & spin waves \\
Regularization: cutoff $\Lambda$ & lattice spacing $a$ \\
Renormalization: $\Lambda \to \infty$ & continuum limit $a \to 0$ \\
Changes in the vacuum & phase transitions \\
\bottomrule
\end{tabular}
\end{table}

Let me start with a quick summary of the Euclidean theory. Consider the generic 2-point correlation function $\Gamma_M(t,k) = \langle 0| \int d^3x\, e^{-ik\cdot x}\, T[F(x,t)F^\dagger(0,0)] |0\rangle$ in Minkowski space. The Heisenberg operator $F(x,t)$ can be written as
\begin{equation}
F(x,t) = e^{iHt-ipx} F e^{-iHt+ipx}
\label{eq:5.3}
\end{equation}
so that
\begin{equation}
\Gamma_M(t,k) = \langle 0| \int d^3x\, e^{-ik\cdot x}\, e^{iHt-ipx} F e^{-iHt+ipx} F^\dagger |0\rangle
\label{eq:5.4}
\end{equation}
where, in dropping the time-ordered product, it is assumed that $t > 0$. Using the invariance of the vacuum under time and space translations and assuming $F^\dagger = F$ for brevity, this becomes
\begin{equation}
\begin{aligned}
\Gamma_M(t,k) &= \langle 0| \int d^3x\, e^{-ik\cdot x}\, F e^{-iHt+ipx} F |0\rangle \\
&= \langle 0| \delta(p-k)\, F e^{-iHt} F |0\rangle .
\end{aligned}
\label{eq:5.5}
\end{equation}

Since the time dependence is explicit, we can Wick rotate to Euclidean space using $t \to -i\tau$ and keeping $H$ unchanged
\begin{equation}
\Gamma_E(t) = \langle 0| \delta(p-k)\, F e^{-H\tau} F |0\rangle .
\label{eq:5.6}
\end{equation}
To get the path integral representation we divide the time interval into $N$ steps $\tau = Na$ and insert a complete set of states $|\phi\rangle\langle\phi|$ at each step.
\begin{equation}
\Gamma_E(t) = \delta(p-k) \sum_{\phi_1 \ldots \phi_N} \langle 0|F|\phi_1\rangle \langle \phi_1| e^{-Ha} |\phi_2\rangle \ldots \ldots \langle \phi_{N-1}| e^{-Ha} |\phi_N\rangle \langle \phi_N| F |0\rangle
\label{eq:5.7}
\end{equation}

The discussion of how, by this introduction of the complete set of states, one trades the operator form of the quantum mechanical amplitude by the sum over paths is standard and thus not repeated. On this ``lattice'' the transfer matrix $T \equiv \exp(-Ha)$ plays the role of the time translation operator, from which the continuum Hamiltonian is obtained as
\begin{equation}
H = \lim_{a\to 0} -\frac{1}{a} \log T .
\label{eq:5.8}
\end{equation}
In order for $H$ to be a self-adjoint Hamiltonian it is necessary that $T$ is a symmetric, bounded, and positive operator acting on a Hilbert space of states with a positive norm. The necessary condition for this, is the Osterwalder-Schrader reflection positivity [24]. That this condition holds for LQCD has been shown by an explicit construction by M.~L\"uscher [25]. Therefore, the correlation functions in the desired physical Minkowski theory can be obtained by analytical continuation of their Euclidean counterparts.

Assuming that the conditions for a sensible Euclidean theory have been established, the vacuum to vacuum amplitude (also called the generating functional or the partition function), in presence of a source $J$, is
\begin{equation}
Z[J] = \int d\phi\, e^{-S_E + J\phi}
\label{eq:5.9}
\end{equation}
where $S_E$ is the Euclidean action. Correlation functions can be obtained from this in the standard way by differentiating $\log Z[J]$ with respect to $J$, for example $\langle 0|\phi(x)\phi(0)|0\rangle = \frac{\delta}{\delta J(x)} \frac{\delta}{\delta J(0)} \log Z[J]\big|_{J=0}$.

Lastly, I summarize the relations between Euclidean and Minkowski quantities. The basic tool is the repeated insertion of a complete set of momentum states in Eq.~\eqref{eq:5.7}. I use the relativistic normalization of states
\begin{equation}
\langle k|p\rangle = 2E L^3\, \delta_{k_x p_x} \delta_{k_y p_y} \delta_{k_z p_z} \xrightarrow{L\to\infty} 2E(2\pi)^3 \delta^3(p-k) .
\label{eq:5.10}
\end{equation}
Then, for a spectrum of stable isolated states $|n\rangle$ with masses $M_n$ that couple to an interpolating field $F$,
\begin{equation}
\Gamma_E(\tau) = \sum_n \frac{\langle 0|F|n\rangle \langle n|F|0\rangle}{2M_n} e^{-M_n \tau} .
\label{eq:5.11}
\end{equation}
Now it is easy to obtain the connection between these $M_n$ and the physical spectrum. A fourier transform in Euclidean time followed by a rotation to Minkowski space gives
\begin{equation}
\int d\tau\, e^{ip_4 \tau}\, \frac{e^{-M_n |\tau|}}{2M_n}
= \frac{1}{2M_n (M_n - ip_4)} + \frac{1}{2M_n (M_n + ip_4)}
= \frac{1}{M_n^2 + p_4^2}
\xrightarrow{p_4 \to -iE} \frac{1}{M_n^2 - E^2} .
\label{eq:5.12}
\end{equation}

The $M_n$ are precisely the location of the poles in the propagator of $|n\rangle$. This is why one can simply short step the necessary analytical continuation as the masses measured from the exponential fall-off are the pole masses of the physical states. For $\tau \to \infty$ only the lowest state contributes due to the exponential suppression, and this provides a way of isolating the lightest mass state.

This simple connection holds also for matrix elements. Consider, for $\tau_2 \gg \tau_1 \gg 0$, the three-point function
\begin{equation}
\begin{aligned}
&\int d^3x\, d^3y\, e^{-ik\cdot x} e^{-iq\cdot y}\, \langle 0| T[F_f(x,\tau_2) O(y,\tau_1) F_i(0,0)] |0\rangle \\
&= \sum_a \sum_b \frac{\langle 0|F_f|a\rangle e^{-M_a (\tau_2 - \tau_1)}}{2M_a}\, \langle a|O|b\rangle\, \frac{\langle b|F_i|0\rangle e^{-M_b \tau_1}}{2M_b} .
\end{aligned}
\label{eq:5.13}
\end{equation}

The two factors on the left and right of $\langle a|O|b\rangle$ are each, up to one decay amplitude, the same as the 2-point function discussed in Eq.~\eqref{eq:5.11}. They, therefore, provide, for large $\tau_2 - \tau_1$ and $\tau_1$, the isolation of the lowest states in the sum over states. These can be gotten from a calculation of the 2-point function. Consequently, $\langle a|O|b\rangle$ measured in Euclidean simulations is the desired matrix element between these states and no analytic continuation to Minkowski space is necessary.

These arguments break down once the sum over states is not just over stable single particle states. Consider, for example, an $F$ that is the $\rho$ meson interpolating field. In that case there are also the $2\pi$, $4\pi$ intermediate states with a cut beginning at $E = 2M_\pi$. The physical characteristics of $\rho$ meson (mass and width) can be represented by a complex pole in the propagator that, in Minkowski space, lies on the second sheet under the cut. This physical pole can no longer be reproduced by measuring the ``energy'' given by the exponential fall-off of Euclidean correlation function at large $\tau$ [26]. It requires measuring the correlation function for all intermediate states and momenta (all $\tau$), and then doing the fourier transform in $\tau$ followed by the analytical continuation $p_4 \to -iE$.

For the same reason there does not exist a simple way for calculating $n$-particle scattering amplitudes. A general 2-particle scattering amplitude (4-point correlation function) is complex precisely because the intermediate states involve a sum over resonances and a continuum with all values of momenta. In principle this can be obtained from Euclidean correlation functions but only by measuring them with essentially infinite precision for all possible $\tau$ and then doing the analytical continuation. In practice this is a hopeless possibility. The one exception is the amplitude at threshold which is real. M.~L\"uscher has shown how the 2-particle scattering amplitudes at threshold can be calculated from the energy shift in a finite volume, the difference between the 2-particle energy and the sum of the two particles in isolation [27--30].
